\subsection{本质有限型代数}

设 \(k\) 为环。设 A 为一个 \(k\) -代数，\(\mathbf{x} = (x_i)_{i \in I}\) 为 A 的元素的一个族；用 \(\mathbf{A}'\) 表示由 \(x_i\) 生成的 A 的子代数。若族 \(\mathbf{x}\) 满足下列条件，则称它为 \(k\) -代数 A 的一个本质生成族：对 A 的每个元素 \(a\)，存在一个元素 \(s\)（属于 \(\mathbf{A}'\) 且在 A 中可逆），使 \(sa\) 属于 \(\mathbf{A}'\)。这等价于说：对每个 \(a \in \mathbf{A}\)，存在多项式 P 和 Q（在 \(k[(\mathrm{X}_i)_{i \in I}]\) 中），使 \(\mathrm{Q}(\mathbf{x})\) 在 A 中可逆且有 \(a = \mathrm{P}(\mathbf{x})\mathrm{Q}(\mathbf{x})^{-1}\) 。

我们将称一个 \(k\) -代数 A 是本质有限型的，如果它具有一个有限本质生成族。这等价于说：存在一个有限型的 \(k\) -代数 A' 和 A' 的一个乘性子集 S，使 \(k\) -代数 A 同构于 \(\mathrm{S}^{-1}\mathrm{A}'\) 。

例.--- 1) 说域 K 的扩张 L 是一个本质有限生成 K-代数，意味着它是 A, V, 第 11 页, 定义 2 意义下的有限生成扩张。K-代数 L 是有限型的，仅当它是 K 的有限次扩张 (V, § 3, 第 4 节, 定理 3 的推论 3)。

\begin{enumerate}
\def\labelenumi{\arabic{enumi})}
\setcounter{enumi}{1}
\tightlist
\item
  为使一个 \(k\) -代数是局部的且本质有限型的，必须且只需它同构于一个 \(k\) -代数（形如 \(\mathsf{A}_{\mathfrak{p}}\)），其中 A 是一个有限生成 \(k\) -代数，\(\mathfrak{p}\) 是 A 的素理想 (参见 II, § 2, 第 5 节, 命题 11 (iii))。
\end{enumerate}

命题 1. - 若环 k 是 Noether 的，则每个本质有限生成 k-代数都是 Noether 环。

这由 III, § 2, 第 10 节, 定理 2 的推论 3 及 II, § 2, 第 4 节, 命题 10 的推论 2 得出。

由定义，我们导出下列性质：

命题 2.--- a) 本质有限型 k-代数的每个商代数都是本质有限型的。

\begin{enumerate}
\def\labelenumi{\alph{enumi})}
\setcounter{enumi}{1}
\item
  本质有限生成 k-代数的每个分式环都是本质有限生成 k-代数。
\item
  有限个本质有限型 \(k\) -代数的乘积 \(k\) -代数是本质有限型的。
\item
  设 \(k \to k'\) 为环同态；对每个本质有限型 \(k\) -代数 A，\(k'\) -代数 \(A_{(k')} = k' \otimes_k \Lambda\) 都是本质有限型的。
\end{enumerate}

推论.- 设 A 为本质有限型 \(k\) -代数，B 为 Noether 的 \(k\) -代数。则环 \(\mathrm{A} \otimes_k \mathrm{B}\) 是 Noether 的。

事实上，它是一个本质有限型 B-代数 (命题 2, d))，于是应用命题 1。

命题 3. 设 \(k\) 为环，A 为一个线性拓扑化的形式平滑 \(k\) -代数，B 为 A-代数。若 A 在 \(k\) 上本质有限型且 B 在 A 上本质有限型，则 B 在 \(k\) 上本质有限型。

设 \(\rho: A \to B\) 为典范映射。设 \(x = (x_i)_{i \in I}\) 为 \(k\) -代数 \(A\) 的一个有限本质生成族，\(A'\) 为它生成的子代数；设 \(y = (y_j)_{j \in J}\) 为 \(A\) -代数 \(B\) 的一个有限本质生成族。设 \(B'\) 为子 \(k\) -代数（\(B\) 中的），由 \(\rho(x_i)\) 与 \(y_j\) 生成。设 \(b \in B\)；按假设存在多项式 \(P\) 和 \(Q\)（在 \(A[(Y_j)_{j \in J}]\) 中），使 \(Q(y)\) 在 \(B\) 中可逆且有 \(Q(y)b = P(y)\) 。\(P\) 和 \(Q\) 的非零系数只有有限多个；存在多项式 \(R \in k[(X_i)_{i \in I}]\)，使 \(R(x)\) 在 \(A\) 中可逆且 \(R(x)P\) 和 \(R(x)Q\) 属于 \(A'[[(Y_j)_{j \in J}]\) 。于是 \(\rho(R(x))Q(y)\) 在 \(B\) 中可逆，并且有

\[
\rho (\mathrm{R} (\mathbf {x})) \mathrm{Q} (\mathbf {y}) b = \rho (\mathrm{R} (\mathbf {x})) \mathrm{P} (\mathbf {y}) \in \mathrm{B} ^ {\prime}.
\]

于是元素 \(\rho(x_i)\)（对 \(i \in I\)）与 \(y_j\)（对 \(j \in J\)）构成 \(k\) -代数 B 的一个本质有限生成族。

推论. -- 两个本质有限型 k-代数的张量积是本质有限型 k-代数。

事实上，设 A 和 B 为两个本质有限生成 k-代数。则 \(A \otimes_{k} B\) 在 A 上本质有限型 (命题 2, d))，从而在 k 上本质有限型 (命题 3)。

